\documentclass[11pt,a4paper]{article}
\usepackage[T1]{fontenc}
\usepackage[utf8]{inputenc}
\usepackage{lmodern}
\usepackage{microtype}
\usepackage{geometry}
\usepackage{listings}
\usepackage{xcolor}
\usepackage{hyperref}
\usepackage{parskip}
\usepackage{booktabs}

\geometry{a4paper, margin=2.5cm}

\definecolor{codebg}{RGB}{245,245,245}
\definecolor{codeframe}{RGB}{200,200,200}
\definecolor{keyword}{RGB}{0,100,180}
\definecolor{comment}{RGB}{100,140,100}
\definecolor{string}{RGB}{180,60,60}

\lstset{
  backgroundcolor=\color{codebg},
  frame=single,
  rulecolor=\color{codeframe},
  basicstyle=\ttfamily\small,
  keywordstyle=\color{keyword}\bfseries,
  commentstyle=\color{comment}\itshape,
  stringstyle=\color{string},
  breaklines=true,
  numbers=left,
  numberstyle=\tiny\color{gray},
  numbersep=6pt,
  showstringspaces=false,
  tabsize=4,
}

\hypersetup{
  colorlinks=true,
  linkcolor=blue,
  urlcolor=blue,
}

\title{\textbf{Wilo-Pumpe Teil 2: Core 1 richtig nutzen}\\[4pt]
  \large MicroPython verwaltet beide Kerne --- so laeuft eigener C-Code trotzdem}
\author{make Magazin $\cdot$ Projekt-Doku}
\date{April 2026 \quad \texttt{github.com/gerontec/wilo\_pwm}}

\begin{document}

\maketitle
\thispagestyle{empty}

\noindent\rule{\textwidth}{0.4pt}
\begin{center}
\textit{Der RP2040 hat zwei Kerne --- aber MicroPython v1.28 nutzt beide.
Wer Core 1 direkt startet, crasht das System.
Dieser Artikel zeigt, wie man trotzdem eigenen C-Code auf Core 1 bekommt:
ueber MicroPythons Thread-System, mit 26 Millionen Polling-Iterationen pro Sekunde.}
\end{center}
\noindent\rule{\textwidth}{0.4pt}

\bigskip

% -----------------------------------------------------------------------
\section{Die Fehlannnahme: Core 1 liegt brach}

Eine verbreitete Annahme beim RP2040 unter MicroPython lautet:
\emph{``Core 0 laeuft MicroPython, Core 1 ist frei.''}
Das stimmt nicht.

MicroPython v1.28 fuer den rp2-Port kompiliert mit
\texttt{MICROPY\_PY\_THREAD=1} --- Threading ist standardmaessig aktiv.
Die Datei \texttt{mpthreadport.c} registriert einen eigenen
\texttt{core1\_entry\_wrapper} und ruft bei Bedarf
\texttt{multicore\_launch\_core1\_with\_stack()} auf.

Der Fehler liegt nahe: Ein eigenes
\texttt{multicore\_launch\_core1(meine\_funktion)} direkt aus C aufzurufen
kollidiert mit dieser Infrastruktur --- mit vorhersehbarem Ergebnis:

\begin{lstlisting}[language=bash, caption={Typisches Fehlerbild beim direkten Core-1-Start}]
# Pico bootet, USB erscheint kurz...
[usb 1-1: device descriptor read/64, error -71]
# ...und verschwindet wieder. Kein ttyACM0, kein WebREPL.
\end{lstlisting}

Der Crash tritt auf bevor MicroPython die Python-VM initialisiert hat,
weil Core 1 sofort in eine enge Polling-Schleife springt und dabei
den USB-Stack stoert.

% -----------------------------------------------------------------------
\section{Die richtige Schnittstelle: mp\_thread\_create}

MicroPython stellt eine offizielle C-API bereit, um Code auf Core 1
zu starten: \texttt{mp\_thread\_create()}.
Diese Funktion leitet den Aufruf durch
\texttt{core1\_entry\_wrapper} --- sie ist kompatibel mit MicroPythons
interner Multicore-Verwaltung.

\begin{lstlisting}[language=C, caption={C-Modul: Core 1 sicher starten}]
#include "py/runtime.h"
#include "py/mpthread.h"
#include "pwmfb_core1_shared.h"
#include "pwmfb_core1_task.h"

static mp_obj_t mp_c1_start(void) {
    static bool started = false;
    if (!started) {
        size_t stack_size = 1024;
        mp_thread_create(pwmfb_core1_task, NULL, &stack_size);
        started = true;
    }
    return mp_const_none;
}
\end{lstlisting}

Der \texttt{static bool}-Guard verhindert, dass
\texttt{mp\_thread\_create} bei jedem Python-\texttt{import} erneut
aufgerufen wird --- wichtig, weil MicroPythons Soft-Reset das Modul
neu initialisiert.

% -----------------------------------------------------------------------
\section{Der Task: Polling mit Batch und Yield}

Ein weiteres Detail: MicroPythons Thread-System erwartet,
dass der Thread gelegentlich Zeit abgibt.
Eine nackte \texttt{while(1)}-Schleife ohne jeden Yield
startet, crasht dann aber den Interpreter.
Die Loesung ist ein Batch-Ansatz:

\begin{lstlisting}[language=C, caption={pwmfb\_core1\_task.c --- Batch-Polling mit sleep\_us}]
#define POLL_BATCH  10000u   // Iterationen pro Runde

void *pwmfb_core1_task(void *arg) {
    (void)arg;
    uint32_t last_t   = TIMER_TIMELR;
    uint8_t  last_lvl = (SIO_GPIO_IN & PIN_MASK) ? 1 : 0;

    while (1) {
        for (uint32_t i = 0; i < POLL_BATCH; i++) {
            g_pwmfb.poll_count++;
            uint8_t lvl = (SIO_GPIO_IN & PIN_MASK) ? 1 : 0;
            if (lvl != last_lvl) {
                uint32_t now  = TIMER_TIMELR;
                uint32_t diff = now - last_t;
                if (diff >= DEBOUNCE_US) {
                    if (lvl == 1) g_pwmfb.low_us  = diff;
                    else          g_pwmfb.high_us = diff;
                    g_pwmfb.last_upd = now;
                    g_pwmfb.ready    = 1;
                    g_pwmfb.edge_count++;
                } else {
                    g_pwmfb.debounce_drop++;
                }
                last_t   = now;
                last_lvl = lvl;
            }
        }
        sleep_us(1);   // kurz yielden -- haelt MicroPython stabil
    }
    return NULL;
}
\end{lstlisting}

10\,000 Polling-Iterationen dauern auf dem RP2040 bei 125\,MHz
etwa 80\,us, dann folgt 1\,us Pause.
Das ergibt in der Praxis \textbf{26 Millionen Polls pro Sekunde}
bei gleichzeitig stabilem MicroPython-Betrieb.

% -----------------------------------------------------------------------
\section{Shared Memory: Keine Locks noetig}

Core 1 schreibt, Core 0 liest.
Die Struct liegt im statischen \texttt{.data}-Segment der Firmware:

\begin{lstlisting}[language=C, caption={pwmfb\_core1\_shared.h}]
typedef struct {
    volatile uint32_t high_us;       // HIGH-Dauer letzter Flanke
    volatile uint32_t low_us;        // LOW-Dauer
    volatile uint32_t last_upd;      // Zeitstempel (TIMER_TIMELR)
    volatile uint32_t poll_count;    // Gesamt-Iterationen
    volatile uint32_t edge_count;    // Erkannte Flanken
    volatile uint32_t debounce_drop; // Verworfene Flanken
    volatile uint32_t debounce_us;   // aktueller Schwellwert (adaptiv)
    volatile uint8_t  ready;
    volatile uint8_t  _pad[3];
} pwmfb_shared_t;

extern pwmfb_shared_t g_pwmfb;
\end{lstlisting}

Auf dem RP2040 sind 32-Bit-Schreibzugriffe auf den SRAM atomar
(Single-Cycle-SRAM, 32-Bit-Bus, ausgerichtete Adressen).
Ein Mutex ist fuer diese Lesemuster nicht noetig.

% -----------------------------------------------------------------------
\section{Adaptiver Debounce: automatisch fuer jede Frequenz}

Ein fester Debounce-Schwellwert funktioniert nur in einem engen
Frequenzbereich.
Bei 75\,Hz (Periode 13,3\,ms) sind 50\,us harmlos ---
bei 10\,kHz (Periode 100\,us) wuerden bereits 50\,\% aller Flanken
verworfen, bei 1\,MHz (Periode 1\,us) jede einzelne.

Die Loesung: Core 1 berechnet den Schwellwert nach jeder erfolgreichen
Messung neu --- 5\,\% der gemessenen Periode, geklemmt zwischen 2\,us
und 5000\,us:

\begin{lstlisting}[language=C, caption={Adaptive Debounce nach jeder Flanke}]
// Nach erfolgreicher Messung (innerhalb des Polling-Batch):
uint32_t period  = g_pwmfb.high_us + g_pwmfb.low_us;
uint32_t new_deb = period / 20;   // 5% der Periode
if (new_deb < DEBOUNCE_US_MIN) new_deb = DEBOUNCE_US_MIN; // min 2 us
if (new_deb > DEBOUNCE_US_MAX) new_deb = DEBOUNCE_US_MAX; // max 5000 us
g_pwmfb.debounce_us = new_deb;
deb = new_deb;   // sofort im laufenden Batch wirksam
\end{lstlisting}

Der Schwellwert passt sich bei jedem Frequenzwechsel automatisch an ---
kein Neustart, kein Konfigurationsparameter.
Der aktuelle Wert ist per \texttt{get\_health()} ablesbar und wird
per MQTT als \texttt{C1\_deb\_us} publiziert:

\begin{center}
\begin{tabular}{lll}
\toprule
\textbf{Signal} & \textbf{Periode} & \textbf{Debounce (adaptiv)} \\
\midrule
75\,Hz    & 13\,333\,us & 666\,us \\
1\,kHz    & 1\,000\,us  & 50\,us \\
10\,kHz   & 100\,us     & 5\,us \\
500\,kHz  & 2\,us       & 2\,us (Minimum) \\
\bottomrule
\end{tabular}
\end{center}

% -----------------------------------------------------------------------
\section{Health-Check: Ist Core 1 noch am Leben?}

Das Python-Interface gibt drei Kenngroessen zurueck:

\begin{lstlisting}[language=Python, caption={Health-Check aus MicroPython}]
import pwmfb_core1
pwmfb_core1.start()

# get_health() -> (poll_count, edge_count, debounce_drop, ready, debounce_us)
poll, edges, drops, ready, deb = pwmfb_core1.get_health()
print(f"Polls/s: ~{poll // uptime_s:,}")
print(f"Flanken total: {edges}")
print(f"Debounce-Verwuerfe: {drops}")
\end{lstlisting}

Typischer Messwert nach 14 Sekunden Laufzeit:

\begin{center}
\begin{tabular}{ll}
\toprule
\textbf{Kenngroesse} & \textbf{Wert} \\
\midrule
poll\_count  & 46\,825\,388 \\
edge\_count  & 2\,777 (bei 75\,Hz $\approx$ 150 Flanken/s) \\
debounce\_drop & 0 \\
debounce\_us  & 666\,us (adaptiv bei 75\,Hz) \\
ready        & 1 \\
\bottomrule
\end{tabular}
\end{center}

\texttt{poll\_count} wachst kontinuierlich --- bleibt er stehen,
ist Core 1 haengengeblieben.
MQTT traegt beide Werte bei jedem 5-Sekunden-Publish aus.

% -----------------------------------------------------------------------
\section{Frequenzmessung: 75 Hz garantiert}

Die Wilo iPWM-Pumpe sendet exakt 75\,Hz.
Core 1 misst HIGH- und LOW-Dauer jeder Periode direkt per
Registerlesezugriff auf \texttt{TIMER\_TIMELR} (1-us-Aufloesung).
Die Python-Seite berechnet daraus Frequenz und Duty-Cycle:

\begin{lstlisting}[language=Python, caption={Frequenzberechnung im Python-Wrapper}]
high_us, low_us, age_us = pwmfb_core1.get_raw()
T    = high_us + low_us      # Periodenlaenge in us
freq = 1_000_000.0 / T       # Frequenz in Hz
duty = high_us / T * 100.0   # Duty-Cycle in Prozent
\end{lstlisting}

Gemessene Werte im Betrieb:
\texttt{75.0\,Hz}, Duty-Cycle je nach Betriebspunkt 29--42\,\%.

% -----------------------------------------------------------------------
\section{Build: USER\_C\_MODULES statt natmod}

Im Gegensatz zum natmod-Ansatz aus Teil\,1 wird das Core-1-Modul
als \texttt{USER\_C\_MODULES} in die Firmware gebaut ---
es benoetigt Zugriff auf MicroPythons interne Thread-API,
die im natmod-Kontext nicht verfuegbar ist.

\begin{lstlisting}[language=bash, caption={Firmware bauen}]
cd ~/micropython/ports/rp2
git checkout v1.28.0
git submodule update --init

make BOARD=RPI_PICO_W \
     USER_C_MODULES=/home/gh/python/pwmfb_core1/micropython.cmake \
     -j$(nproc)
# -> build-RPI_PICO_W/firmware.uf2
\end{lstlisting}

Das MicroPython-Dateisystem (main.py, .mpy-Dateien) bleibt
unveraendert. Nach dem Flash einfach per WebREPL oder mpremote
neu deployen.

\begin{center}
\begin{tabular}{lll}
\toprule
\textbf{Ansatz} & \textbf{Max.\,Freq.} & \textbf{Besonderheit} \\
\midrule
Python IRQ (pure)   & $\sim$1--2\,kHz  & kein Build \\
C-Natmod (.mpy)     & $\sim$10\,kHz    & kein Build, portabel \\
Core 1 via Thread   & $\sim$500\,kHz   & Custom Firmware noetig \\
PIO (Assembler)     & $\sim$62\,MHz    & Custom Firmware + ASM \\
\bottomrule
\end{tabular}
\end{center}

Fuer 75\,Hz ist jeder Ansatz ausreichend.
Der Wert liegt im Debugging: \texttt{poll\_count} und
\texttt{edge\_count} per MQTT machen Core-1-Ausfall sofort sichtbar.

\bigskip
\noindent\rule{\textwidth}{0.4pt}
\textbf{Code, Firmware \& Deploy-Script:}
\url{https://github.com/gerontec/wilo_pwm}

\end{document}
